% ═══════════════════════════════════════════════════════════════════
%  PAGE DE GARDE — format EMSI
% ═══════════════════════════════════════════════════════════════════
%
%  Ce fichier est appelé de DEUX endroits :
%    · main.tex             — le rapport complet ;
%    · couverture-seule.tex — un document d'une page, pour régler la
%                             couverture sans recompiler les 102 pages.
%
%  Il n'existe donc qu'une seule version de la couverture : ce qu'on ajuste
%  ici se retrouve des deux côtés.
%
%  Les valeurs affichées (titre, noms, filière, année) ne sont PAS ici : elles
%  vivent en tête de main.tex et de couverture-seule.tex, dans le bloc encadré
%  de ▼▼▼ / ▲▲▲.
%
%  ⚠️ Deux passes de compilation sont nécessaires : `remember picture` doit
%  connaître la position des coins de la page, ce qu'il n'apprend qu'au
%  passage suivant. Une couverture décalée au premier essai est normale.
%
%  Toutes les positions sont mesurées depuis le COIN SUPÉRIEUR GAUCHE de la
%  page (`current page.north west`), en centimètres, y négatif vers le bas.
%  Pour descendre un bloc, augmenter son deuxième nombre.
%
%  Repères verticaux : logos −2,3 · filière −6,6 · cadre du titre −7,9 à −12,9
%  · « Réalisé par » −16,5 · « Encadré par » −20,6 · année −27,6.
\begin{titlepage}
\thispagestyle{empty}
\begin{tikzpicture}[remember picture, overlay]

  % ---- Bande jaune en haut à gauche ----
  \fill[jauneEMSI] ($(current page.north west)+(0.4cm,-1.1cm)$)
    rectangle ($(current page.north west)+(8cm,-1.5cm)$);

  % ---- Logos : l'école à gauche, le projet à droite ----
  \node[anchor=north west] at ($(current page.north west)+(0.7cm,-2.3cm)$) {%
    \logoCouverture{EMSI.png}{2.1cm}
  };
  \node[anchor=north east] at ($(current page.north east)+(-1cm,-2.3cm)$) {%
    \logoCouverture{logo.png}{2.3cm}
  };

  % ---- Bande jaune sous les logos ----
  \fill[jauneEMSI] ($(current page.north west)+(3.6cm,-5.3cm)$)
    rectangle ($(current page.north east)+(-0.4cm,-5.7cm)$);

  % ---- Nature du travail + filière ----
  \node[anchor=north east] at ($(current page.north east)+(-1cm,-6.6cm)$) {%
    \begin{minipage}{9cm}\raggedleft
      {\large\bfseries \nouvelleNature}\\[4pt]
      {\normalsize\bfseries\color{vertEMSI} \nouvelleFiliere}
    \end{minipage}
  };

  % ---- Ombre + cadre arrondi du titre, avec ses deux attaches jaunes ----
  \fill[ombre, rounded corners=8pt] ($(current page.north west)+(0.9cm,-8.1cm)$)
    rectangle ($(current page.north west)+(11.2cm,-13.1cm)$);
  \draw[jauneEMSI, line width=3pt, rounded corners=8pt]
    ($(current page.north west)+(0.7cm,-7.9cm)$)
    rectangle ($(current page.north west)+(11cm,-12.9cm)$);
  \fill[jauneEMSI] ($(current page.north west)+(11cm,-8.3cm)$)
    rectangle ($(current page.north east)+(-0.4cm,-8.7cm)$);
  \fill[jauneEMSI] ($(current page.north west)+(11cm,-12.1cm)$)
    rectangle ($(current page.north west)+(11.4cm,-12.9cm)$);
  % La césure est désactivée ici seulement : un titre coupé en « ap-plication »
  % au milieu d'une couverture se voit tout de suite. Le reste du rapport, lui,
  % garde la césure — sans elle, le texte justifié se troue.
  %
  % \hyphenchar et non \hyphenpenalty : dans un nœud TikZ, la pénalité est
  % rétablie avant que le paragraphe ne soit coupé, donc sans effet. Retirer le
  % caractère de césure de la police, après l'avoir sélectionnée, tient.
  \node[anchor=center, align=center, text width=9.6cm]
    at ($(current page.north west)+(5.85cm,-10.4cm)$) {%
    \normalsize\bfseries\selectfont
    \hyphenchar\font=-1\relax
    \nouvelleTitre
  };

  % ---- Filet jaune de séparation ----
  \fill[jauneEMSI] ($(current page.north west)+(1cm,-13.9cm)$)
    rectangle ($(current page.north east)+(-1cm,-14.05cm)$);

  % ---- Bloc « Réalisé par » ----
  % La barre est calée sur UN auteur. Pour un binôme, ajouter la seconde ligne
  % dans le minipage ci-dessous et descendre le −17.9 à −18.5.
  \fill[jauneEMSI] ($(current page.north west)+(1.1cm,-16.5cm)$)
    rectangle ($(current page.north west)+(1.4cm,-17.9cm)$);
  \node[anchor=north west] at ($(current page.north west)+(1.7cm,-16.6cm)$) {%
    \begin{minipage}{12cm}
      {\bfseries\color{vertEMSI} Réalisé par~:}\\[6pt]
      \nouvelleAuteur
    \end{minipage}
  };

  % ---- Bloc « Encadré par » ----
  \fill[jauneEMSI] ($(current page.north west)+(1.1cm,-20.6cm)$)
    rectangle ($(current page.north west)+(1.4cm,-23.4cm)$);
  \node[anchor=north west] at ($(current page.north west)+(1.7cm,-20.7cm)$) {%
    \begin{minipage}{15.5cm}
      {\bfseries\color{vertEMSI} Encadré par~:}\\[6pt]
      \nouvelleEncadrantPedagogique, Encadrant Académique,
      EMSI \nouvelleVille.\\[6pt]
      \nouvelleEncadrantEntreprise, Encadrant Professionnel,
      \nouvelleTitreCourt, \nouvelleVille.
    \end{minipage}
  };

  % ---- Filet jaune de bas de page ----
  \fill[jauneEMSI] ($(current page.north west)+(1cm,-26.7cm)$)
    rectangle ($(current page.north east)+(-1cm,-26.85cm)$);

  % ---- Année universitaire ----
  \node[anchor=north] at ($(current page.north)+(0cm,-27.6cm)$) {%
    {\large\bfseries Année Universitaire~: \nouvelleAnnee}
  };

\end{tikzpicture}
\end{titlepage}
